% ---------------------------------------------------------------------
% Argomento: legge di Ohm, potenza dissipata
% Livello: base   |   Tempo stimato: 6 min
% ---------------------------------------------------------------------
\begin{esercizio}
\dato{R}{470}{680}{560}
\dato{V}{9}{12}{15}
\calcola[I]{I}{\x{V}/\x{R}}
\calcola[P]{P}{\x{V}^2/\x{R}}

Un resistore da \mis{R}{\ohm} è collegato a una tensione di \mis{V}{\volt}.
\begin{parts}
  \part[2] Calcola la corrente che lo attraversa.
  \part[2] Calcola la potenza dissipata.
\end{parts}

\begin{solution}
  $I = V/R \approx \mis{I}{\ampere}$;\qquad
  $P = V^2/R \approx \mis{P}{\watt}$
  \criterio{1 pt formula, 1 pt risultato; $-0{,}5$ se manca o è sbagliata l'unità.}
\end{solution}
\end{esercizio}
